%% kapa: exam=2 round=28 year=2017 period=3 subject=compensation_law subject_ko=감정평가 및 보상법규 seq=3 number=3 points=20
%% source: Q-Net 감정평가사 기출문제 (cst00309) artlSeq=5201712 file=raw/1008/5201712/2017년 감정평가사 2차시험 기출문제(감정평가 및 보상법규).pdf
%% conversion: kordoc 4.24.51 -> tools/gen2.py (auto); review=proofread-sonnet; flags=none
\begin{longquestion}{【문제 3】}
지목은 대(垈)이지만 그 현황이 인근 주민의 통행에 제공된 사실상 도로인 토지를 대상으로 「도시 및 주거환경정비법」에 따른 매도청구권을 행사하는 경우와 「공익사업을 위한 토지 등의 취득 및 보상에 관한 법률」에 따른 수용재결이 행하여지는 경우에 관하여 다음 물음에 답하시오. (20점) \par
\medskip\noindent 1) 매도청구권 행사에 따른 쟁송절차와 수용재결에 따른 보상금을 다투는 쟁송절차의 차이점을 설명하시오. (10점) \par
\medskip\noindent 2) 토지의 감정평가방법과 그 기준에 있어 매도청구권이 행사되는 경우와 수용재결이 행하여지는 경우의 차이점을 설명하시오. (10점) \par
\end{longquestion}
